\begin{problem}[Warmup]{(a)}
    Consider the vectors $\mathbf{v} = (2, -1)$ and $\mathbf{w} = (2, 3)$ in $\mathbb{R}^2$. Compute the following, and draw diagrams to interpret these computations geometrically: \[\mathbf{v} + \mathbf{w}, \quad \mathbf{v} - \mathbf{w}, \quad 2\mathbf{w}, \quad \tfrac{1}{2}\mathbf{v}, \quad -\mathbf{v}.\]
\end{problem}
\unsolved

\begin{problem}[Warmup]{(b)}
    Explain the difference between $0$ and $\mathbf{0}$.
\end{problem}
\unsolved

\begin{problem}[Warmup]{(c)}
    Let $V = \{f \in \mathbb{Z}_3[x] : \deg(f) \le 2\}$; explain why this is a vector space over $\mathbb{Z}_3$. List all vectors and all scalars in this case.
\end{problem}
\unsolved

\begin{problem}[Warmup]{(d)}
    Let $V = \{a + b\pi : a, b \in \mathbb{Q}\}$. Is this a vector space over $\mathbb{Q}$? Over $\mathbb{R}$?
\end{problem}
\unsolved

\begin{problem}[Warmup]{(e)}
    Let $V$ be a vector space over $\mathbb{R}$, and let $\mathbf{v} \in V$. Explain what $-\mathbf{v}$ and $(-1)\mathbf{v}$ mean, according to the definitions of our notation; why are they equal?
\end{problem}
\unsolved

\begin{problem}[Warmup]{(f)}
    Is $\mathbb{C}$ a vector space over $\mathbb{Q}$? Is $\mathbb{Q}$ a vector space over $\mathbb{C}$?
\end{problem}
\unsolved

\begin{problem}[Warmup]{(g)}
    Is $\mathbb{Z}[x]$ a vector space over $\mathbb{Z}$?
\end{problem}
\unsolved

\begin{problem}{1}
    Check Example 40.4: That is, prove that $\mathbb{Q}[x]$ is a vector space over
    $\mathbb{Q}$.
\end{problem}
\unsolved

\begin{problem}{2}
    Check Example 40.5: That is, prove that $\mathbb{Q}_n[x]$ is a vector space over
    $\mathbb{Q}$, although it is not a subring of $\mathbb{Q}[x]$.
\end{problem}
\unsolved

\begin{problem}{3}
    Check Example 40.3: That is, let $F$ be a field, and let $F^n$ be the set of
    $n$-tuples with entries from $F$. Prove that $F^n$ is a vector space over $F$.
\end{problem}
\unsolved

\begin{problem}{4}
    Let $V$ be a vector space over the field $F$. Suppose that $r, s \in F$ and
    $\mathbf{0} \neq \mathbf{v} \in V$. Prove that if $r\mathbf{v} = s\mathbf{v}$, then
    $r = s$.
\end{problem}
\unsolved

\begin{problem}{5}
    Let $V$ be a vector space over the field $F$. Suppose that $0 \neq r \in F$. Define
    the function
    \[\varphi_r : V \to V\]
    by $\varphi_r(\mathbf{v}) = r\mathbf{v}$. Prove that $\varphi_r$ is a one-to-one onto
    function that preserves addition. That is, $\varphi_r$ is an \emph{additive group
    isomorphism} from $(V, +)$ onto itself.
\end{problem}
\unsolved

\begin{problem}{6}
    Let $V$ be a vector space over the field $F$. Suppose that $r, s \in F$ and
    $\mathbf{v} \in V$. Prove that $r(s\mathbf{v}) = s(r\mathbf{v})$.
\end{problem}
\unsolved

\begin{problem}{7}
    Show that $M_2(\mathbb{R})$, the two-by-two matrices with entries from $\mathbb{R}$, is
    a vector space over $\mathbb{R}$.
\end{problem}
\unsolved

\begin{problem}{8}
    Show that $M_{m,n}(F)$, the $m$-by-$n$ matrices with entries from the field $F$, is a
    vector space over $F$.
\end{problem}
\unsolved

\begin{problem}{9}
    If $F$ and $E$ are fields with $F \subseteq E$, show that $E$ is a vector space over
    $F$. This is an important example of a vector space in subsequent chapters.
\end{problem}
\unsolved

\begin{problem}{10}
    Let
    $\mathbb{Q}(\sqrt[3]{2}) = \{a + b\sqrt[3]{2} + c\sqrt[3]{4} : a, b, c \in \mathbb{Q}\}$.
    (In Exercise 38.5 you showed this is a field.) Show that $\mathbb{Q}(\sqrt[3]{2})$ is a
    vector space over $\mathbb{Q}$.
\end{problem}
\unsolved

\begin{problem}{11}
    Let $V$ be the set of real-valued functions with addition defined by
    \[(f + g)x = f(x) + g(x).\]
    For $c \in \mathbb{R}$ and $f \in V$, the scalar multiple of $f$ by $c$ is
    $cf(x) = c(f(x))$. Show $V$ is a vector space over $\mathbb{R}$.
\end{problem}
\unsolved

\begin{problem}{12}
    Prove that $\mathbb{Z}$ is not a vector space over $\mathbb{Z}_p$, where $p$ is a
    positive prime integer.
\end{problem}
\unsolved

\begin{problem}{13}
    Prove that $\mathbb{Z}$ is not a vector space over $\mathbb{Q}$.
\end{problem}
\unsolved
